% 1장 서론. 출처: docs/framing.md 1절 · docs/proposal.md

\section{서론}

\subsection{동기와 연구 질문}

토지거래허가제는 일정 면적을 넘는 토지 거래에 관할 구청장의 허가를 요구하는
제도다. 가격에 상한을 두지도, 세금을 매기지도 않는다. 거래 자체에 \textbf{허가
절차라는 마찰}을 붙일 뿐이다. 서울시는 2020년 6월 23일 국제교류복합지구 인근
네 개 법정동(대치·삼성·청담·잠실)을 지정했고, 2021년 4월 압구정·여의도·목동·성수로
넓혔으며, 2025년 2월 대부분을 해제했다가 39일 만에 더 넓게 재지정했다.

이 제도가 거래를 줄였다는 것은 새로운 발견이 아니다. 2020년 6월 지정을 다룬
국내 선행연구 세 편이 모두 ``거래 건수가 60\% 넘게 줄었다''고 이미 적었다
\citep{jeon2025, lee2025, yang2025}. 거래에 허가를 붙였으니 거래가 줄어드는 것은
제도의 정의에 가깝다. 그것만 확인하는 작업이라면 의미가 적다.

\textbf{이 연구가 묻는 것은 그다음이다.} 다섯 가지로 나눈다.

\begin{enumerate}[leftmargin=2em, itemsep=2pt]
  \item \textbf{얼마나} 줄었나 --- 선행연구의 60\%는 날짜를 달력 월로 세어
        지정 전 거래가 섞여 희석된 값일 수 있다.
  \item \textbf{언제까지} 갔나 --- 일시 충격인지 시장이 계속 얼어붙었는지에
        따라 정책의 비용이 다르다.
  \item \textbf{누구에게} 달랐나 --- 규제 설계의 어느 부분이 실제로 작동했는지.
  \item \textbf{어디로} 번졌나 --- 규제 밖 단지로 수요가 옮겨 갔다면 규제는
        문제를 푼 것이 아니라 옮긴 것이다.
  \item \textbf{가격은} 잡혔나 --- 토허제의 명분이 가격 안정이기 때문이다.
\end{enumerate}

이 다섯 질문에 대한 답을 미리 한 문장으로 요약하면 이렇다.
\textbf{토허제의 명분은 가격 안정인데 실제로 작동한 것은 거래 동결이었다. 거래는
확실하게 끊었지만 가격을 잡았다는 증거는 약하다.} 그리고 그 거래 동결조차
모두에게 같지 않았다 --- 허가 대상에서 빠진 소형 아파트는 오히려 거래가 늘었다.

\subsection{기여 사항}

\textbf{첫째, 거래 단위 미시 패널을 지었다.} 선행연구 세 편은 모두 동-월 집계
자료를 쓴다. 이 연구는 실거래 ``상세'' 자료로 거래 한 건마다 계약일, 전용면적,
층, 단지 식별자를 붙이고, 정책 공고 원문에서 읽은 지정 구간을 법정동별로 결합해
거래 시점의 규제 상태를 붙였다. 잘게 쪼갠 것 자체가 기여는 아니다. 기여는
\textbf{미시 자료가 아니면 보이지 않는 것을 보였다}는 데 있다(3장).

\textbf{둘째, 설계 선택이 결과를 가르는 정도를 정량화했다.} 합리적인 설계 조합
192가지를 전부 추정해, 선행연구끼리 결론이 엇갈린 이유가 ``지역을 어떻게
나눴느냐''(2.3\,\%p)가 아니라 ``사건 시점을 어디에 맞췄느냐''(39.6\,\%p)임을
보였다(2.4절).

\textbf{셋째, 효과의 이질성이 허가 대상의 면적 기준에서 왔음을 자연실험으로
갈랐다.} 처음에는 대출 규제가 원인이라고 보았으나 경계를 옮겨 보니 규제 문턱인
15억원에 계단이 없었다. 2022년 6월 허가 기준이 대지 18\,m$^2$에서 6\,m$^2$로
강화된 사건이 두 가설을 갈랐다(6장).

\textbf{넷째, 그 답으로 세운 예측이 다른 자료에서 확인됐다.} ``2025년 재지정은
면적 기준이 6\,m$^2$이므로 소형도 함께 줄어야 한다''고 미리 적었고, 실제로 소형
거래가 71.5\% 줄었다(6.8절). 그리고 2021년 4월 지정이라는 별개 사건에서 같은
크기의 효과가 나왔다(9.6절).

\textbf{다섯째, 전 과정을 공개하고 클린 환경에서 재현을 확인했다.} 원자료 수집부터
그림과 표 생성까지 \pipeSteps 단계 파이프라인 하나로 돈다. 문서에 적힌 수치가 분석 결과와
어긋나지 않는지, 아무도 만들지 않는 산출물이 남아 있지 않은지, 의존 목록이 실제
import와 맞는지를 매 실행마다 점검한다(3.5절).

\subsection{본 보고서의 구성}

2장은 제도와 선행연구를 다루고, 선행연구 세 편을 이 연구의 자료로 직접 재현한다.
그리고 세 편의 결론이 왜 엇갈렸는지를 설계 선택의 분해로 답한다. 3장은 자료와
시스템이다. 4장은 식별 전략, 5장부터 8장까지가 다섯 질문에 대한 답이며, 9장은
강건성, 10장은 논의와 한계, 11장이 결론이다.

설계를 바꾼 판단과 그 근거는 결정기록 31편에 남겼다(부록 A). 본문에서 ``0014''와
같이 가리키는 것은 그 번호다.
